\chapter{結果與討論}

\section{研究結果}

請依照研究問題或實驗設計呈現主要結果。結果段落應先描述觀察到的現象，再解釋其意義，避免將推論與原始觀察混在一起。

\begin{table}[htbp]
  \centering
  \caption{範例表格}
  \label{tab:example}
  \begin{tabular}{lll}
    \toprule
    項目 & 說明 & 備註 \\
    \midrule
    A & 第一項結果 & 請替換為正式內容 \\
    B & 第二項結果 & 請替換為正式內容 \\
    \bottomrule
  \end{tabular}
\end{table}

\section{討論}

請根據研究目的與文獻缺口討論結果。可說明結果是否支持研究假設、與既有研究的異同、可能原因與實務意涵。

\section{效度與限制}

請討論研究限制，例如資料代表性、方法假設、實驗條件、外部效度、量測誤差或可能的偏誤來源。

\section{本章小結}

本章呈現研究結果並討論其意義、限制與可能影響。下一章將總結本論文的主要貢獻與未來方向。
